\documentclass{article}
\usepackage[UTF8]{ctex}
\usepackage{amsmath,amssymb,amsfonts}
\usepackage{enumitem}
\usepackage{geometry}
\geometry{a4paper, margin=1in}
% === 错题本排版优化（不参与题目标识） ===
\setlength{\parindent}{0pt}
\setlength{\parskip}{0.35em plus 0.10em minus 0.08em}
\setlength{\emergencystretch}{3em}
\setlist{itemsep=0.25em, topsep=0.40em, parsep=0pt}
\allowdisplaybreaks[2]
% === 错题本排版优化结束 ===
\begin{document}
\section*{中国科学技术大学2026年高等代数真题}
\section*{99．中国科学技术大学2026年线性代数与解析几何试卷}
一．填空题。\\

\subsection{题目1}
1．直线 \(\left\{\begin{array}{l}6 x-2 y+3 z=9, \\ 3 x-3 y+z=6 .\end{array}\right.\) 求点 \((0,0,0)\) 到该直线的距离 \(\_\_\_\_\)．\\
\(\_\_\_\_\)

\subsection{题目2}
2．直线 \(\left\{\begin{array}{l}2 x+y-z=2, \\ 3 x-y+z=3 .\end{array}\right.\) 求该直线绕 \(y\) 轴旋转所得旋转曲面的方程 \(\_\_\_\_\)．

\subsection{题目3}
3．设 \(|A|=\left|\begin{array}{cccc}1 & 2 & 1 & -7 \\ 2 & 7 & 2 & 4 \\ 1 & 1 & 0 & -2 \\ 17 & 16 & 19 & 11\end{array}\right|\) ，求 \(A_{41}+A_{42}+A_{43}+A_{44}=\)

\subsection{题目4}
4．已知 \(A\) 的伴随矩阵 \(A^{*}=\left(\begin{array}{lll}1 & 1 & 1 \\ 1 & 2 & 1 \\ 1 & 3 & 3\end{array}\right)\) ，求原矩阵 \(A=\) \(\_\_\_\_\)

\subsection{题目5}
5．在三次多项式空间 \(P[x]_{3}\) 中，求多项式 \(1+x+x^{2}+x^{3}\) 在基 \(\left\{2,1+x,-x+2 x^{2}, x+x^{3}\right\}\) 下的坐标\\
\(\_\_\_\_\)．

\subsection{题目6}
6．二次型 \(2 x_{1}^{2}+2 x_{2}^{2}+2 x_{3}^{2}-2 x_{1} x_{2}-2 x_{1} x_{3}-2 x_{2} x_{3}\) 的规范型为 \(\_\_\_\_\)．

\subsection{题目7}
7．矩阵 \(\left(\begin{array}{rrr}13 & 16 & 16 \\ -5 & -7 & -6 \\ -6 & -8 & -7\end{array}\right)\) 的若尔当标准型为 \(\_\_\_\_\)．\\
\(\_\_\_\_\)

\subsection{题目8}
8．设 \(M=\left(\begin{array}{cc}1 & 0 \\ 1 & -1\end{array}\right)\) ，线性变换 \(\mathcal{A}\) 满足 \(\mathcal{A} x=M x\) ，求 \(\mathcal{A}\) 的特征值及重数 \(\_\_\_\_\)．

\subsection{题目9}
9．设 \(A=\left(\begin{array}{ccc}2026 & 2 & -2 \\ 2 & 2029 & -4 \\ -2 & -4 & 2029\end{array}\right)\) ，记 \(M=\left\{B \in M_{3}(\mathbb{R}) \mid A B=B A\right\}\)（即所有与 \(A\) 可交换的 3 阶实矩阵构成的集合），求 \(\operatorname{dim} M=\) \(\_\_\_\_\)．

\subsection{题目10}
10．在欧氏空间中，向量 \(d_{1}=(1,1,1,-1), d_{2}=(2,0,-2,0)\) 张成的子空间为 \(V\) ，求向量 \(\beta=\) \((1,-1,1,0)\) 在 \(V\) 上的投影向量 \(\_\_\_\_\)．

二．解答题。\\

\subsection{题目11}
1．设矩阵

\[
A=\left(\begin{array}{cccc}
1+a & 1 & \cdots & 1 \\
2 & 2+a & \cdots & 2 \\
\vdots & \vdots & & \vdots \\
n & n & \cdots & n+a
\end{array}\right)
\]

求 \(\operatorname{det} A\) ，判断 \(a\) 为何值时，\(A X=0\) 有非零解，并求该非零解．

\subsection{题目12}
2．设矩阵 \(A=\left(\begin{array}{ccc}1 & 2 & -3 \\ -1 & 4 & -3 \\ 1 & a & 5\end{array}\right)\) 有二重特征值，判断 \(a\) 为何值时 \(A\) 可对角化，并求可逆矩阵 \(P\) ，使得\(P^{-1} A P\) 为对角阵．

\subsection{题目13}
3．设矩阵 \(A=\left(\begin{array}{lll}1 & 0 & 0 \\ 2 & 2 & 0 \\ 3 & 3 & 3\end{array}\right)\) ．

\begin{itemize}
  \item[（1）] 求上三角矩阵 \(B\) 与正交矩阵 \(Q\) ，使得 \(A=B Q\) ．
  \item[（2）] 判断 \(B, Q\) 是否唯一？若不唯一则有几对？并论证理由。
\end{itemize}

\subsection{题目14}
4．设 \(A\) 为 \(n\) 阶可逆实方阵，且 \(\operatorname{tr}(A)=n\) ，证明：存在可逆矩阵 \(P\) ，使得 \(P^{-1} A P\) 是对角元全为1的矩阵。

\subsection{题目15}
5．设 \(A\) 为 \(n\) 阶可逆复矩阵，记 \(A=\left(\begin{array}{cc}A_{11} & A_{12} \\ A_{21} & A_{22}\end{array}\right)\) ，其中 \(A_{11}\) 为 \(m \times(n-m)\) 型矩阵 \((1<m<n)\) ，若 \(A_{21}\) 在 \(A\) 中所有到代数余子式全为 0 ，证明：\(n \leq 2 m\) 。

\subsection{题目16}
6．设 \(A\) 为元素全为整数的对称方阵，且对任意非零行向量 \(z\)（其分量均为非负整数），总有 \(z A z^{\mathrm{T}} \geq 0\) 。

\begin{itemize}
  \item[（1）] 证明：对任意非零行向量 \(z\)（其分量均为非负实数），总有 \(z A z^{\mathrm{T}} \geq 0\) 。
  \item[（2）] 判断：是否存在非零行向量 \(z\)（其分量均为非负整数），使得 \(z A z^{\mathrm{T}}=0\) 。
\end{itemize}
\end{document}
